%% title: Hypothesis Testing - Size, Power, Neyman-Pearson and the Three Testing Principles
\documentclass[11pt]{article}
\usepackage{metricsstyle}

\begin{document}

\handoutheader{MLE III}{Some notes on hypothesis testing: size and power, consistency, the Neyman--Pearson lemma, and the Wald, likelihood ratio and Lagrange multiplier principles}{}

\runin{Previously} \emph{MLE~I--II}: score $g(\theta)=\partial\log L/\partial\theta$, Hessian $H(\theta)$,
information matrix $\mathcal I$, and $\sqrt n(\hat\theta-\theta_{0})\overset{a}{\sim}N(0,\mathcal I^{-1})$.

\section{Hypothesis testing}

Hypothesis testing is concerned with whether the data appear to favour or disfavour a particular
description of nature.

There is one particular hypothesis we are interested in, the \emph{null hypothesis} $H_{0}$. If the
data fall into a particular region of the sample space, called the \emph{critical region}, the test
is said to \emph{reject} $H_{0}$; otherwise it \emph{accepts}. In testing there are only two
outcomes (reject, accept), whereas in estimation there is a continuum of possible outcomes.

\topic{Errors, size and power}
Decision making centres around the following possibilities:
\begin{center}
\begin{tabular}{l|c|c}
 & $H_{0}$ true & $H_{0}$ false\\ \hline
 Accept & correct & error II\\
 Reject & error I & correct
\end{tabular}
\end{center}
\begin{align*}
 \alpha&=\Pr(\text{Type I error})\quad\text{(the \emph{size} of the test)},\\
 \beta&=\Pr(\text{Type II error}),\\
 1-\beta&=\text{power of the test}.
\end{align*}
The power gives the probability of \emph{correctly} rejecting $H_{0}$.

In comparing tests, the standard notion of optimality is based on size and power. Within a class of
tests, one is said to be \emph{best} if it has the maximum power among all tests with size less than
or equal to some $\alpha$ (e.g.\ $\alpha=0.05$).

To make such conditions operational it is necessary to specify how the data are generated when
$H_{0}$ is false. This is the role of the \emph{alternative hypothesis}. By specifying an
alternative, the critical region can be tailored to look for deviations from $H_{0}$ in the
direction of the alternative.

\topic{Set-up}
Let $y$ be a $T\times1$ random vector drawn from the joint density $f(y,\theta)$, where $\theta$ is a
$k\times1$ vector of unknown parameters and $\theta\in\Theta$, the parameter space. Then
\begin{align*}
 H_{0}&:\ \theta\in\Theta_{0}\subset\Theta,\\
 H_{A}&:\ \theta\in\Theta_{A}\subset\Theta,
\end{align*}
with $\Theta_{0}\cap\Theta_{A}=\varnothing$. Frequently $\Theta_{A}=\Theta-\Theta_{0}$, i.e.\ $\Theta_{A}$ is
the complement of $\Theta_{0}$ with respect to $\Theta$.

For a critical region $C_{T}$, the size $\alpha_{T}$ is given by
\[ \alpha_{T}=\Pr\{\,y\in C_{T}\mid\theta\in\Theta_{0}\,\}, \]
and the power is
\[ \pi_{T}(\theta)=\Pr\{\,y\in C_{T}\mid\theta\in\Theta_{A}\,\}. \]

\topic{Definition (consistency)}
A test is \emph{consistent} if
\[ \lim_{T\to\infty}\pi_{T}(\theta)=\pi(\theta)=1,\qquad\theta\in\Theta_{A}, \]
i.e.\ a consistent test will always reject $H_{0}$ asymptotically if it is false. In other words,
Type II errors are eliminated from large samples.

\section{The Neyman--Pearson lemma}

Limited in scope, but useful because it provides intuition. Suppose that $H_{0}$ and $H_{A}$ each
consist of a single value of $\theta$:
\[ H_{0}:\ \theta=\theta_{0},\qquad H_{A}:\ \theta=\theta_{1}. \]
Let $L(\theta_{0})$ and $L(\theta_{1})$ be the likelihood under $H_{0}$ and $H_{1}$ respectively.

\topic{Lemma}
If there exists a critical region $C$ of size $\alpha$ and a constant $k$ such that
\[ \frac{L(\theta_{0})}{L(\theta_{1})}\le k\ \text{ inside }C
   \qquad\text{and}\qquad
   \frac{L(\theta_{0})}{L(\theta_{1})}>k\ \text{ outside }C, \]
then $C$ is the \emph{most powerful} critical region (Type II errors are minimised among all
regions of size $\alpha$).

\emph{Note.} If $H_{0}$ is true, we would expect $L(\theta_{0})$ to be large relative to $L(\theta_{1})$,
so the ratio of the likelihoods would tend to fall outside the critical region. If $H_{A}$ is true,
$L(\theta_{0})/L(\theta_{1})$ would tend to be small, $\le k$, inside $C$.

\topic{The likelihood ratio}
\[ \lambda=\frac{L(\theta_{0})}{L(\theta_{1})} \]
is the \emph{likelihood ratio}, and
\[ -2\log\lambda=2\bigl[\log L(\theta_{1})-\log L(\theta_{0})\bigr]\ \overset{a}{\sim}\ \chi^{2}_{r}, \]
where $r$ is the number of restrictions (the distribution is asymptotic, under $H_{0}$, and for the
case where $\theta_{1}$ is the unrestricted maximum likelihood estimate).

\section{Three principles of testing a hypothesis}

\begin{figure}[h]
\centering
\begin{tikzpicture}[scale=1.05]
 \draw[ecaxis] (0,0) -- (7,0) node[below] {$\theta$};
 \draw[ecaxis] (0,0) -- (0,4.4) node[left] {$\log L$};
 \draw[line width=0.9pt,domain=0.4:6.4,samples=60,smooth] plot (\x,{3.7-(\x-5)^2*0.13});
 % restricted point theta_0 = 2.6 : f = 3.7-0.13*5.76 = 2.9512, slope = 0.26*(5-2.6)=0.624
 \draw[dashed] (0,3.7) node[left] {$\log L(\theta_{1})$} -- (5,3.7);
 \draw[dashed] (0,2.9512) node[left] {$\log L(\theta_{0})$} -- (2.6,2.9512);
 \draw[dashed] (2.6,0) -- (2.6,2.9512);
 \draw[dashed] (5,0) -- (5,3.7);
 \draw[thin] (1.2,{2.9512+0.624*(1.2-2.6)}) -- (4.2,{2.9512+0.624*(4.2-2.6)}) node[above right] {$g(\theta_{0})$};
 \node[ecnode] at (6.3,3.0) {$g(\theta_{1})=0$};
 \draw[|-|] (-1.55,2.9512) -- (-1.55,3.7);
 \node[ecnode] at (-1.95,3.33) {LR};
 \draw[|-|] (2.6,-0.3) -- (5,-0.3);
 \node[ecnode] at (3.8,-0.62) {W};
 \node[ecnode] at (2.6,-0.65) {$\theta_{0}$};
 \node[ecnode] at (5,-0.65) {$\theta_{1}$};
\end{tikzpicture}
\end{figure}

\begin{flatnum}
 \item \textbf{Wald} principle: based on $(\hat\theta_{1}-\theta_{0})$.
 \item \textbf{Likelihood ratio} principle: uses $\log\lambda=\log L(\hat\theta_{0})-\log L(\hat\theta_{1})$.
 \item \textbf{Lagrange multiplier} principle: based on $g(\theta_{0})$, where
       $g(\theta)=\partial\log L(\theta)/\partial\theta$.
\end{flatnum}

Here $\theta_{1}$ denotes the unrestricted parameter (estimate). Then
\begin{align*}
 \text{W looks at}\quad&\lVert\theta_{0}-\theta_{1}\rVert,\\
 \text{LM looks at}\quad&\lVert g(\theta_{0})-g(\theta_{1})\rVert\equiv\lVert g(\theta_{0})\rVert,\\
 \text{LR looks at}\quad&\lVert\log L(\theta_{1})-\log L(\theta_{0})\rVert,
\end{align*}
where $\lVert\cdot\rVert$ denotes distance. The distance function gives you the \emph{units of
measurement} that the arguments in the function are apart: for W, how many units $\theta_{0}$ is away
from $\theta_{1}$; for LM, how many units $g(\theta_{0})$ is from $g(\theta_{1})$. The identity for LM holds
because $g(\theta_{1})=0$ at the unrestricted maximum.

\emph{Q: what is this function?} The next section answers it: the units are set by the information matrix.

\section{The distance function: the three statistics}

Recall from the handout on maximum likelihood (\emph{MLE~I}). Let
$L(\theta\mid y)=\prod_{t=1}^{n}f(y_{t}\mid\theta)$ be the likelihood, where $f(y_{t}\mid\theta)$ is the density
of observation $t$. Then
\[ \log L(\theta\mid y)=\ell(\theta\mid y)=\sum_{t=1}^{n}\log f(y_{t}\mid\theta), \]
so $\log L$ is a sum of random variables. The gradient and Hessian are
\[ g(\theta)=\frac{\partial\log L}{\partial\theta}=\sum_{t=1}^{n}\frac{\partial\log f(y_{t}\mid\theta)}{\partial\theta},
   \qquad
   H(\theta)=\frac{\partial^{2}\log L}{\partial\theta\,\partial\theta'}
   =\sum_{t=1}^{n}\Bigl\{\frac{\partial^{2}\log f(y_{t}\mid\theta)}{\partial\theta_{i}\,\partial\theta_{j}}\Bigr\}, \]
the braces holding the typical $ij$ element. From the first-order condition $g(\hat\theta)=0$ we derived that
\[ \frac1{\sqrt n}\,g(\theta_{0})\overset{a}{\sim}N(0,\mathcal I),
   \qquad
   \sqrt n(\hat\theta-\theta_{0})\overset{a}{\sim}N\bigl(0,\mathcal H^{-1}\mathcal I\,\mathcal H^{-1}\bigr), \]
with $\mathcal I$ the information matrix and $\mathcal H=\operatorname{plim}H/n$. Also
\[ \mathcal H^{-1}\mathcal I\,\mathcal H^{-1}=-\mathcal H^{-1}=\mathcal I^{-1}, \]
because if the $y_{t}$ are independent, $\mathcal H=-\mathcal I$.

\topic{LM}
Therefore $\lVert g(\theta_{0})-g(\theta_{1})\rVert=\lVert g(\theta_{0})\rVert$ becomes
\[ \frac1n\,g(\theta_{0})'\,\mathcal I(\theta_{0})^{-1}g(\theta_{0})
   =-\frac1n\,g(\theta_{0})'\,\mathcal H(\theta_{0})^{-1}g(\theta_{0}). \]
One replaces $\mathcal H$ by $H$ and estimates under $H_{0}$: replace $\theta_{0}$ by $\tilde\theta$, the
``restricted'' estimator, and let $\hat\theta$ be the ``unrestricted'' estimator. Hence
\[ LM=\frac1n\,g(\tilde\theta)'\bigl[-H(\tilde\theta)/n\bigr]^{-1}g(\tilde\theta)
   =g(\tilde\theta)'\bigl[-H(\tilde\theta)\bigr]^{-1}g(\tilde\theta)\ \overset{a}{\sim}\ \chi^{2}_{r}. \]

\topic{Wald}
\[ W=n(\hat\theta-\theta_{0})'\{\widehat{\operatorname{var}}\}^{-1}(\hat\theta-\theta_{0})
   =n(\hat\theta-\theta_{0})'\bigl[-H(\hat\theta)/n\bigr](\hat\theta-\theta_{0})\ \overset{a}{\sim}\ \chi^{2}_{r}, \]
where $\widehat{\operatorname{var}}=\bigl[-H(\hat\theta)/n\bigr]^{-1}$ estimates the asymptotic covariance matrix of
$\sqrt n(\hat\theta-\theta_{0})$.

\topic{LR}
And finally
\[ LR=2\bigl[\log L(\hat\theta)-\log L(\tilde\theta)\bigr]\ \overset{a}{\sim}\ \chi^{2}_{r}. \]

\clearpage
\begin{transcriptnotes}
 \item Source: \emph{MLE3.pdf}, nine photographed pages of the handwritten handout headed
       ``(III) Some notes on Hypothesis Testing'', paginated 1--9 (page 9 holds only the LR line). The
       text is restated and the derivations expanded; the section headings are \textbf{added}. No
       equation numbers in the handout. No date. Transposes ($^{T}$) are primes; the script
       $\mathcal H$ (limiting Hessian) is distinguished from the sample Hessian $H$.
 \item Pages 1--2: the handout's table has $H_{0}$ true/false across and accept/reject down; ``error~II'' is
       accepting a false $H_{0}$ and ``error~I'' rejecting a true one. The set-up is written
       $H_{0}\colon\theta\in\Theta_{0}\subset\Theta$ and $H_{A}\colon\theta\in\Theta_{A}\subset\Theta$ with empty intersection.
 \item Pages 3--4: ``$\theta=\theta_{1}$'' (simple alternative) and the likelihood ratio
       $\lambda=L(\theta_{0})/L(\theta_{1})$ are as written. \textbf{Added:} the qualification that
       $-2\log\lambda\sim\chi^{2}_{r}$ is asymptotic and holds when $\theta_{1}$ is the unrestricted
       maximum likelihood estimate (the lemma itself is for two simple hypotheses).
 \item Page 5: the figure is redrawn: $\log L$ against $\theta$, with the tangent $g(\theta_{0})$ at
       $\theta_{0}$, the maximum at $\theta_{1}$ where $g(\theta_{1})=0$, and the brackets LR (vertical) and W
       (horizontal). The handout writes $\log\lambda=\log L(\hat\theta_{0})-\log L(\hat\theta_{1})$; kept.
 \item Page 6: ends with the question ``What is this function?''; the answer is given on pages 7--8
       (the information matrix) and is made explicit here. The handout writes
       $\lVert g(\theta_{0})-g(\theta_{1})\rVert\equiv\lVert g(\theta_{0})\rVert$ without the reason; the
       reason ($g(\theta_{1})=0$) is \textbf{added}.
 \item Page 8: \textbf{Corrected:} the handout's LM statistic carries the factor $1/n$ together with the
       sample Hessian $H$ ($=\sum_{t}$, as defined on page 7), and the Wald statistic's second line drops
       the factor $n$ ($(\hat\theta-\theta_{0})'[(-H/n)^{-1}]^{-1}(\hat\theta-\theta_{0})$). With the
       definitions of page 7 the consistent forms are $LM=g'(-H)^{-1}g$ (equivalently
       $\frac1n g'(-H/n)^{-1}g$) and $W=n(\hat\theta-\theta_{0})'(-H/n)(\hat\theta-\theta_{0})$; these are
       written out above. ``$\operatorname{var}(\hat\theta)$'' in $W$ is read as the covariance matrix of
       $\sqrt n(\hat\theta-\theta_{0})$. Degrees of freedom $r$ (number of restrictions) as in the handout.
 \item Original page map: p.\,1 hypothesis testing, outcomes, error table, size/power; p.\,2 power, best
       test, role of alternative, set-up $\Theta_{0}$, $\Theta_{A}$; p.\,3 size, power, consistency,
       Neyman--Pearson set-up; p.\,4 lemma, likelihood ratio, $-2\log\lambda\sim\chi^{2}_{r}$; p.\,5 three
       principles and figure; p.\,6 distances, units of measurement; p.\,7 recall of ML results;
       p.\,8 LM and Wald statistics; p.\,9 LR statistic.
\end{transcriptnotes}

\end{document}
